\documentclass{docs}
\begin{document}
    \begin{titlepage}
        \begin{center}
            \vspace*{1cm}
            
            \Huge
            \textbf{The KScr Language}
            
            \vspace{0.5cm}
            \LARGE
            Definition of the KScr Language
            
            \vspace{1.5cm}
            
            \textbf{comroid}
            
            \vspace{0.8cm}
            
            \includegraphics[width=0.4\textwidth]{../img/kscr-icon.png}
            
            \vfill
            
            The KScr Language as defined by this document is a compiled high-level language whose goal is to provide runtime optimizations from a different approach.
            
            \vspace{1.2cm}
            
            \Large
            \underline{Team:} \\
            Kaleidox
        \end{center}
    \end{titlepage}
    \fancyhead[C]{{KScr Language}}
    \tableofcontents
    \pagebreak
    \section{Source Model}
    \subsection{Modifiers}
    \subsubsection{Accessibility Modifiers}
    \begin{center}
    \begin{tabular}{c|c}
        Modifier & Function \\
        \hline
        \texttt{public} & Accessible from everywhere \\
        \texttt{internal} & Accessible during compilation; compiled to \texttt{private} \\
        \texttt{protected} & Accessible from inheritors \\
        \texttt{private} & Accessible from inside only
    \end{tabular}
    \end{center}
    \subsubsection{Other Modifiers}
    \begin{center}
    \begin{tabular}{c|c}
        Modifier & Function \\
        \hline
        \texttt{static} & Cannot be invoked dynamically \\
        \texttt{final} & Cannot be overridden or changed \\
        \texttt{abstract} & Must be implemented in inheritors \\
        \texttt{synchronized} & Invocations are synchronized \\
        \texttt{native} & Must be implemented by a native module
    \end{tabular}
    \end{center}
    \subsection{Classes}
    \subsubsection{Class Types}
    A source class can be of different types:
    \paragraph{\texttt{class}-Type}
    A normal class that can be instantiated.
    Allowed modifiers are:
    \begin{center}
    \begin{tabular}{c|c}
        Modifier & Function \\
        \hline
        \texttt{static} & Cannot be instantiated and behaves like a Singleton \\
        \texttt{final} & Cannot be inherited by other classes \\
        \texttt{abstract} & Cannot be instantiated directly and must be inherited by other classes
    \end{tabular}
    \end{center}
    \paragraph{\texttt{enum}-Type}
    An enumeration of runtime-constants that follow a class-like pattern.
    This type does not allow modifiers.
    \paragraph{\texttt{interface}-Type}
    An interface that declares basic structure requirements for implementing classes.
    Cannot be instantiated directly.
    \paragraph{\texttt{annotation}-Type}
    \subsubsection{Type Generics}
    \subsection{Class Members}
    \subsubsection{Static Initializer}
    \subsubsection{Constructors}
    \subsubsection{Methods}
    \subsubsection{Properties}
    \pagebreak
    \section{The KScr VM}
    \subsection{Built-in Types}
    \subsubsection{\texttt{interface void}}
    \label{typeVoid}
    The universal base type. Is implemented by all built-in types and can be implicitly cast to everything.
    \subsubsection{\texttt{class num<T>}}
    The base type of all numerics. Contains numeric subtypes:
    \begin{enumerate}
        \item \texttt{byte} - Type-alias for \texttt{int<8>} % todo: this is not true
        \item \texttt{short} - Type-alias for \texttt{int<16>}
        \item \texttt{int<n = 32>}
        \item \texttt{long} - Type-alias for \texttt{int<64>}
        \item \texttt{float}
        \item \texttt{double}
    \end{enumerate}
    All subtypes can be used directly, or using the Type Generic T, for example: \texttt{int<24> == num<int<24>>}
    \subsubsection{\texttt{class str}}
    \subsubsection{\texttt{class object}}
    \subsubsection{\texttt{class array<T>}}
    \subsubsection{\texttt{class tuple<T...>}}
    \subsubsection{\texttt{class enum<T>}}
    \subsubsection{\texttt{class type<T>}}
    \subsubsection{\texttt{class pipe<T...>}}
    \subsection{Literals}
    \subsubsection{\texttt{null}}
    The \texttt{null}-Literal. Is always of \hyperref[typeVoid]{type \texttt{void}}.
    \subsubsection{\texttt{Numeric Literals}}
    \subsubsection{\texttt{String Literals}}
    A string literal is pre- and superceded by a double-quote \texttt{"} symbol.
    An escaped double-quote \texttt{\textbackslash"} can be contained in the string.
    \paragraph{Planned Behaviour: Interpolation}
    A string supports interpolation using accolades with Formatter-support with the following syntax:
    \begin{verbatim}
        int hex = 1 << 3;
        stdio <<- "hex: {hex:X}"
        // prints "hex: 0x4"
    \end{verbatim}
    \subsubsection{\texttt{Array Literals}}
    Unimplemented.
    \subsubsection{\texttt{Tuple Literals}}
    Unimplemented.
    \subsubsection{\texttt{Other Literals}}
    \paragraph{\texttt{stdio}}
    Constantly represents the program's standard IO stream. \\
    The held value is of type \texttt{pipe<str>}
\end{document}